\subsection{OBS006: the padded multiplicity at the fourth four-row stretch}
\label{obs:fourth-ray-padded}
\paragraph{Statement and scope.}
For the padded permanent pair $(m,n)=(8,9)$, let
$\lambda_d=(9d-28,12,8,8)$ be a partition.
Then $q_{\lambda_d}=a_{\lambda_d}$ in every admissible degree $d$.
In particular, at $d=20$,
\[
 \lambda_{20}=(152,12,8,8),\qquad q_{\lambda_{20}}=a_{\lambda_{20}}=52137.
\]
This is the fourth stretch of $(38,3,2,2)$. Its determinant multiplicity
$K_9(\lambda_{20})$ remains undetermined.

\paragraph{Centered coordinates.}
Write $E=\mathbb C^3$. On the chart $F_0=1,F_1=0$, use normalized polars
$C_j\in\operatorname{Sym}^j E$, so that
$F(x_0,v)=x_0^9+\sum_{j=2}^9\binom9j x_0^{9-j}C_j(v)$.
The localized highest-weight construction in OBS005 identifies every
highest-weight polynomial of the indicated tail with a nonzero polynomial
$I(C_2,\ldots,C_9)$ of tail weight $(12,8,8)$ and weighted degree $28$.
The identification is injective in every fixed outer degree.
For arbitrary lower polars and a linear form $\ell$, actual padded forms
have
\[
 C_9=\ell^9+\sum_{j=2}^8(-1)^j\binom9j\ell^{9-j}C_j.
 \tag{OBS006.1}
\]
Indeed the corresponding ordinary coefficient polynomial vanishes at
$x_0=-\ell$ and hence has the linear factor $x_0+\ell$.
The verified fullness of quaternary $\mathrm{per}_8$ identifies these
forms with points of the required four-variable padded image.

\paragraph{Only a quadratic top-polar term needs attention.}
A term of degree three in $C_9$ would have coefficient of weighted degree
$28-27=1$. There is no generator of weight one after centering, so that term
is absent; higher top-polar degrees are impossible. Write
\[
 I=I_0+I_1[C_9]+A[C_9,C_9].
\]
Here $A$ has lower-polar weighted degree $10$. Its part depending on $C_8$
is necessarily bilinear in $(C_2,C_8)$, since $10-8=2$ and $2\cdot8>10$.
Thus $A=A_0+L(C_2,C_8)$, with $A_0$ independent of $C_8$.
The quadratic nonic representation decomposes as
\[
 \operatorname{Sym}^2(\operatorname{Sym}^9 E)
 =\bigoplus_{k=0,2,4,6,8}S_{(18-k,k)}E.
\]
Only $S_{(12,6)}E$ and $S_{(10,8)}E$ can contribute to first tail row $12$.

\paragraph{The finite leading-coefficient map.}
Isolate the coefficient of $C_8$-degree two and $\ell$-degree ten after
substitution (OBS006.1). It is, up to a nonzero scalar, the map
\[
 \tau\otimes{\rm id}_{\operatorname{Sym}^2 E},\qquad
 \tau:\operatorname{Sym}^2(\operatorname{Sym}^9 E)\otimes
              \operatorname{Sym}^8 E
       \longrightarrow
       \operatorname{Sym}^{10}E\otimes
              \operatorname{Sym}^2(\operatorname{Sym}^8 E),
\]
applied to $L$.
No term of $I_0$ or $I_1[C_9]$ has $\ell$-degree ten.
The contribution of $A_0$ involving two copies of $C_8$ has
$\ell$-degree two, so it does not contribute either.
By the horizontal-strip Pieri rule, an intermediate partition $\nu$
which can yield $(12,8,8)$ on tensoring with $\operatorname{Sym}^2 E$
must satisfy $\nu_2=8$ and $\nu_1+\nu_3=18$.
The complete list is
\[
 (10,8,8),\quad (11,8,7),\quad (12,8,6).
\]
The first two have source multiplicity one and nonzero image under
$\tau$, as explicitly checked in QC010.
The last has source multiplicity two, once from each of
$S_{(10,8)}E\otimes\operatorname{Sym}^8 E$ and
$S_{(12,6)}E\otimes\operatorname{Sym}^8 E$.

Here is an explicit full-rank check for this multiplicity-two part.
Use vector letters $a,b,c$ of degrees $(9,9,8)$, put
$\Delta=[abc]$, $H=[ab]$, $J=[ac]$, $K=[bc]$, and use subscripts $0,1,2$
for their coordinates. Two symmetric-in-$a,b$ highest-weight symbols of
weight $(12,8,6)$ are
\[
 M_1=\Delta^6H^2a_0b_0c_0^2,\qquad
 M_2=\Delta^6JK\,a_0^2b_0^2.
\]
Set $a=e_0$, differentiate once in $b_0$, and average under interchange
of $b,c$. This represents the leading map, up to fixed nonzero
normalization constants. With
\[
 D=b_1c_2-b_2c_1,\quad
 S=b_1^2c_0^2+c_1^2b_0^2,\quad
 P=b_0b_1c_0c_1,
\]
the images are
\[
 M_1\longmapsto \tfrac12D^6S,\qquad
 M_2\longmapsto D^6(\tfrac32S-2P).
\]
They are linearly independent, since $S$ and $P$ are independent
polynomials. The two source symbols therefore form the full
multiplicity-two basis, and $\tau$ is injective on every relevant
intermediate isotypic part. Tensoring this injection with
$\operatorname{Sym}^2 E$ preserves injectivity. Consequently the
vanishing of the isolated coefficient implies $L=0$.

\paragraph{Finishing the whole-space injection.}
Now $A=A_0$ is independent of $C_8$. Isolate the coefficient linear in
$C_8$ and of $\ell$-degree ten. It is the map
\[
 \sigma:\operatorname{Sym}^2(\operatorname{Sym}^9 E)
       \longrightarrow \operatorname{Sym}^{10}E\otimes
                         \operatorname{Sym}^8 E,\qquad
 Q\longmapsto Q(\ell^9,\ell C_8).
\]
It is injective on the entire source: every even-$k$ summand
$S_{(18-k,k)}E$, $0\le k\le8$, occurs once in the target by Pieri,
and its $k$-bracket highest symbol has nonzero image at
$\ell=e_0,\ C_8=e_0^{8-k}e_1^k$. Distinct irreducible summands have
independent images. Thus $A=0$.
The coefficient of $\ell$-degree nine now gives $I_1[\ell^9]=0$.
Pure ninth powers span $\operatorname{Sym}^9 E$, hence $I_1=0$;
then $I_0=0$. Therefore a highest-weight polynomial vanishing on the
padded image is zero, proving injectivity on the full multiplicity
space and the claimed equality $q=a$.

\paragraph{Checks.}
The complete proof above was independently read by both
\texttt{/root/support\_and\_lmr} and
\texttt{/root/quaternary\_certificates}; both reproduced the two
displayed images and the complete Pieri inventory.
The ambient value $52137$ was computed independently by exact
BigInt coefficient dynamic programming with all 24 terms of the
$\mathrm{GL}_4$ Weyl alternation. The same implementation recovered
the certified values $804,1250,2030,2386$.
These checks do not constitute a determinant-rank computation.
